\section{From Harness to Trainable Scaffold}
\label{sec:scaffold}

The \S\ref{sec:harness} evaluation establishes a \emph{family} of active-perception scaffolds, not a single winning harness: at the Qwen3.5-27B backbone the two strongest members, RLM and CodeAct, attain the highest accuracy and are statistically indistinguishable (39.4 and 39.5 ANLS, $n{=}8$), sharing a tool set, prompt body, and VLM-tool perception. The win is the family. The supervised study of \S\ref{sec:training} trains one member, the append-only CodeAct; the on-policy methods left to future work --- policy-gradient reinforcement learning and per-token distillation --- single it out for a sharper reason. The two harnesses differ in how each turn's context is formed across turns, and that difference --- not accuracy --- is what determines whether those on-policy objectives are even well-defined. This section makes that argument: it is why the append-only member is the scaffold to train inside, and the analysis applies whether the training has happened yet or, as for the on-policy methods here, lies ahead.

\subsection{The append-only prefix invariant}

Policy-gradient multi-turn reinforcement learning (Group Relative Policy Optimization, GRPO; Proximal Policy Optimization, PPO; GSPO) and per-token distillation losses are built on an append-only token-prefix invariant: $\text{prompt}_{t+1} = \text{prompt}_t + \text{completion}_t + \text{observation}_t$. Under this invariant the whole rollout is one monotonically growing token sequence --- the context at turn $t$ is an exact prefix of the context at turn $t{+}1$ --- so the trainer recovers every per-step log-probability from a single masked forward pass over the trajectory, scoring each action in exactly the context that generated it \citep{renderers}.

CodeAct preserves this invariant. It maintains an append-only chat transcript --- a system message followed by alternating user and assistant turns --- in which each turn only appends and no earlier token is rewritten. RLM does not. It surfaces a sidecar of metadata describing the variables currently live in the REPL, holds their full values off-context in the interpreter, and re-renders this managed view of state at each turn rather than appending to it. Because the visible context is rebuilt every step, earlier positions change from turn to turn and consecutive turns do not stand in a prefix relation.

Breaking the invariant breaks the policy gradient. Rewriting history makes the observation distribution policy-dependent and non-stationary, violating the assumptions of policy-gradient methods \citep{foldact}, and the rollout is no longer a single growing sequence, so per-step likelihoods must be recovered from separate forward passes over distinct contexts \citep{resum}. Importance ratios and cross-turn credit assignment are then no longer defined over one coherent sequence. Methods that let an agent compress its own history during reinforcement learning report exactly this and repair it with purpose-built machinery: FoldAct stabilizes summary actions with custom losses, and ReSum segments the trajectory at each compression and broadcasts a trajectory-level advantage across segments. An append-only transcript is the structure these objectives are built for; a context-manipulating one needs bespoke repair.

\subsection{Prefix preservation is free}

Selecting the append-only member concedes no accuracy where it is measured. With the corrected CodeAct loop, CodeAct ties RLM at both scales evaluated under matched perception --- 39.5 against 39.4 at the 27B backbone, and 22.3 against 21.1 at the 4B reasoner that the training in \S\ref{sec:training} actually targets ($n{=}8$, both differences within trial-to-trial variance). The structural property that makes CodeAct trainable thus comes at no measured accuracy cost relative to RLM's managed, re-rendered context --- including at the small scale where it matters for training. (On some bases RLM remains the more robust of the two, and the remaining cross-model CodeAct cells await the same corrected re-run.) The conclusion the bridge carries forward is the family-level one --- train inside the append-only member of the winning family, paying nothing in accuracy for a trajectory that established training machinery accepts directly.

\subsection{An accessible ceiling versus a higher one}

The two scaffolds differ in how context length scales. RLM's managed view stays compact across arbitrarily many turns, decoupling the number of reasoning steps from the context window, whereas CodeAct's append-only transcript grows monotonically and can approach the window limit on long documents with many perception calls. In principle this gives RLM a higher ceiling on long-horizon tasks. That the advantage does not surface as higher accuracy is consistent with a specific account: exploiting a compact, re-rendered context requires a policy trained to operate on it, but contemporary models are post-trained almost entirely on append-only conversations, leaving a frozen model off-distribution in such a harness. Managed-context systems indeed tend to match or beat append-only baselines only after regime-specific training --- an untrained context curator underperforms full-context prompting and overtakes it only once trained with reinforcement learning \citep{contextcurator}, and stable context folding requires purpose-built losses \citep{foldact}.

This account is a hypothesis, not a demonstrated result, and its magnitude should not be overstated: frozen models sometimes benefit from compaction with no retraining \citep{resum}, and that benefit shrinks as model capability grows. The conservative conclusion is that an append-only scaffold delivers its full potential under established training, whereas a context-managing scaffold may have a higher but currently less accessible ceiling, reachable only through new and not-yet-established methods. We adopt the prefix-preserving scaffold, for which mature training machinery applies directly; training a policy natively for a context-managing harness is left to future work. The cost of this choice is the append-only transcript's monotonic growth, which on long documents can press against the context window --- a limitation for the longest documents that this preliminary study does not attempt to quantify.
